\section{Candidate Divine Calendar: A Fail-Closed Construction}
\label{sec:divine-calendar}

\subsection{What is now attested}

A striking but carefully delimited numerical fact emerges from first-millennium
Mesopotamian learned spellings.  Six successive multiples of ten can each write
the name of a major deity:
\[
\begin{array}{c|cccccc}
N & 10 & 20 & 30 & 40 & 50 & 60\\ \hline
\text{deity} & \text{Adad} & \text{Shamash} & \text{Sin} & \text{Ea} & \text{Enlil} & \text{Anu}
\end{array}
\]
ORACC independently documents these numerical writings.  This is a genuine
historical number--deity system.  It is \emph{not}, by itself, a calendar.
The guard case is Ištar, whose numerical spelling is 15; the pantheon is
therefore not an exclusively decadal sequence.

\subsection{Six equal active sectors}

The candidate creation calendar is defined on an abstract accumulated phase
(or winding count) $W$, without yet identifying its physical oscillator:
\[
D_n=\{W:10(n-1)\leq W<10n\},\qquad n=1,\ldots,6.
\]
The corresponding endpoints are
\[
B_n=10n\in\{10,20,30,40,50,60\}.
\]
The coincidence between these six endpoints and the six attested divine
numerals is recorded as a \emph{structural candidate}.  Historical dependence,
calendar function, and astronomical meaning remain separate questions.

\subsection{The 60 boundary and the seventh state}

At the endpoint of the sixth active sector,
\[
W=60.
\]
For a phase coordinate only, we may write
\[
\phi_{60}(W)=W\bmod 60,
\qquad
\phi_{60}(60)=0.
\]
This expresses mathematical closure of a base-60 phase.  It must not be
misreported as evidence that ancient scribes possessed a theological concept
of ``zero at Sabbath.''  The historical fact is weaker: sexagesimal counting
uses 60 as a radix transition, and learned divine writing associates 60 with
Anu.

We therefore define the seventh position as a discrete post-closure state,
not as an automatically assumed interval $[60,70)$:
\[
\mathcal K(W)=
\begin{cases}
(\mathrm{ACTIVE},\lfloor W/10\rfloor+1), & 0\le W<60,\\
(\mathrm{POST\mbox{-}CLOSURE},7), & W=60.
\end{cases}
\]
No rule is currently assigned for $W>60$.  This is intentionally fail-closed.

The literary structure of Genesis is compatible with this distinction in one
specific sense.  Days one through six receive the recurring evening--morning
closure, including the sixth day in Gen.~1:31; Gen.~2:2--3 instead describes
the seventh by completion and cessation and does not append another
``evening and morning'' formula.  This observation supports treating day seven
as textually distinct, but does not establish the phase model.

\subsection{The 360 level}

The number 360 is independently well attested in several ancient systems of
time reckoning.  Egypt used a civil structure of twelve 30-day months, each
month divided into three 10-day weeks, followed by five epagomenal days.  In
ancient Judah, Jonathan Ben-Dov argues for a schematic 360-day administrative
and literary year; the priestly flood chronology is particularly diagnostic
because five schematic months correspond to exactly 150 days.

Thus the identity
\[
360=12\times30=36\times10=6\times60
\]
is historically surrounded by real calendrical structures, although the last
factorisation, $6\times60$, is our mathematical re-expression rather than an
attested ancient definition of a ``divine year.''  Later Indian cosmic-time
traditions provide an independent control in which a divine year is explicitly
constructed from 360 divine-day units.

This motivates a provisional hierarchy
\[
10\;\longrightarrow\;60\;\longrightarrow\;360,
\]
with epistemic labels
\[
\text{sector candidate}\;\longrightarrow\;
\text{closure candidate}\;\longrightarrow\;
\text{schematic-year candidate}.
\]
None of these labels fixes a conversion into terrestrial years.

\subsection{Physical clock still unresolved}

The calendar must not be converted into gigayears until its oscillator is
specified.  The correct generic quantity is an integrated phase,
\[
W_X(t)=\frac{1}{2\pi}\int_{t_0}^{t}\Omega_X(\tau)\,d\tau,
\]
not the backward extrapolation of the present Solar orbital period.  A physical
proposal for $X$ must be fixed before fitting chronology, and must survive
uncertainty in the evolving Galactic potential.

\subsection{Falsification conditions}

The divine-calendar candidate is weakened or rejected if (i) the endpoint
sequence requires post-hoc changes, (ii) the physical oscillator is selected
only after seeing the desired age, (iii) the divine numerals prove unrelated
to any temporal or ordered structure beyond scribal shorthand, or (iv) the
$W=60$ boundary does not remain stable under a physically reconstructed
phase history.  Conversely, the model gains support only from predictions
fixed before new evidence is inspected.
